\documentclass[11pt,twocolumn]{article}

% -------------------------------------------------------
% PACKAGES
% -------------------------------------------------------
\usepackage[margin=1in]{geometry}
\usepackage{microtype}
\usepackage{graphicx}
\usepackage{booktabs}
\usepackage{hyperref}
\usepackage{url}
\usepackage{enumitem}
\usepackage{titlesec}
\usepackage{abstract}
\usepackage{fancyhdr}
\usepackage{amsmath}
\usepackage{amssymb}
\usepackage{times}
\usepackage{xcolor}
\usepackage{caption}
\usepackage{subcaption}
\usepackage{array}
\usepackage{parskip}

% -------------------------------------------------------
% PAGE STYLE
% -------------------------------------------------------
\pagestyle{fancy}
\fancyhf{}
\rhead{\small CS437/CS5317/EE414/EE513 --- Deep Learning, Spring 2026}
\lhead{\small SOA Survey}
\rfoot{\small \thepage}
\renewcommand{\headrulewidth}{0.4pt}

% -------------------------------------------------------
% TITLE FORMATTING
% -------------------------------------------------------
\titleformat{\section}
  {\large\bfseries}{\thesection.}{0.5em}{}[\titlerule]
\titleformat{\subsection}
  {\normalsize\bfseries}{\thesubsection.}{0.5em}{}

% -------------------------------------------------------
% ABSTRACT FORMATTING
% -------------------------------------------------------
\renewcommand{\abstractnamefont}{\normalfont\bfseries}
\renewcommand{\abstracttextfont}{\normalfont\small}
\setlength{\absleftindent}{0pt}
\setlength{\absrightindent}{0pt}

% -------------------------------------------------------
% HYPERLINK STYLING
% -------------------------------------------------------
\hypersetup{
    colorlinks=true,
    linkcolor=black,
    citecolor=black,
    urlcolor=blue
}

% -------------------------------------------------------
% TITLE BLOCK
% -------------------------------------------------------
\title{
    \vspace{-1cm}
    {\large \textbf{SOA Survey}} \\[0.3em]
    {\LARGE Early Warning System for Forest Fire} \\[0.5em]
    {\normalsize CS437/CS5317/EE414/EE513 --- Deep Learning, Spring 2026}
}

\author{
    Nausherwan Khan Arzani (27100384) \quad Saad Jamshaid Khan (27100159)
}

\date{}

% -------------------------------------------------------
\begin{document}

\maketitle
\thispagestyle{fancy}

% -------------------------------------------------------
\begin{abstract}
Forest fires pose one of the most destructive natural hazards globally,
threatening ecosystems, wildlife, and human settlements every year.
Early detection before a fire grows uncontrollable is critical, yet
traditional approaches relying on watchtowers, satellite imagery, and
manual inspection suffer from fundamental limitations in spatial
coverage, detection latency, and operational cost. This survey reviews
the state of the art in deep learning-based forest fire and smoke
detection, synthesizing fifteen papers spanning image-based detection,
video analysis, lightweight model design, dataset construction, and
vision-language integration. We examine core architectures including
YOLO variants, Faster R-CNN, attention-based CNNs, and hybrid
transformer models. Our review identifies three central gaps: most
high-accuracy models are too computationally demanding for real-time
deployment on edge hardware such as the Raspberry Pi or Jetson Nano;
the persistent confusion between early-stage smoke and environmental
visual confounders such as clouds, mist, and steam remains unsolved
in fixed-camera deployments; and nighttime and low-light fire detection
is severely underrepresented in both datasets and published methods.
We propose investigating a lightweight detection model augmented with
CLIP-based features and trained with synthetic nighttime augmentation,
evaluated on a custom dataset annotated from WWF-deployed PTZ cameras.
\end{abstract}

% -------------------------------------------------------
\section{Problem Definition}
\label{sec:problem}

Forest fires are among the most severe and rapidly escalating natural
disasters worldwide. Hundreds of millions of hectares of forests are
devastated annually, causing irreversible ecological damage, significant
loss of human life, and enormous economic costs~\cite{martinez2008}.
Studies indicate that approximately 85\% of wildfires in the United
States between 1992 and 2015 were caused by human activity, underscoring
that effective early-warning infrastructure can prevent the majority of
destructive fires~\cite{bouguettaya2022}. Once a fire ignites, the
window for intervention narrows drastically: a small ignition on a
high-risk day can grow out of control within minutes, making early
detection of the first signs of smoke the single most impactful
intervention point available.

This project addresses the problem of
automated fire and smoke detection using footage from Pan-Tilt-Zoom
(PTZ) cameras deployed at fixed vantage points across forest sites.
The task can be formulated in one of three ways depending on the
desired level of output: (1) \textit{binary classification}, where
each image or video frame is labelled as fire/smoke present or absent;
(2) \textit{bounding box detection}, where the spatial location of
fire or smoke within the frame is also predicted; or (3)
\textit{instance segmentation}, where a pixel-level mask is produced.
The system must operate in real time and be deployable on edge hardware
such as a Raspberry Pi or NVIDIA Jetson device, since the cameras are
located at remote forest sites with limited connectivity.

What makes this problem uniquely challenging from a deep learning
perspective is a combination of factors. First, early-stage smoke is
visually ambiguous --- it occupies a small fraction of image pixels,
has poorly defined edges, and its appearance (transparency, colour,
texture) overlaps with benign phenomena such as clouds, haze, and
dust~\cite{wildfire2023}. Second, PTZ cameras are mounted at fixed
elevated positions looking across the horizon, meaning that clouds,
morning mist, and distant smoke appear in the same region of the frame
at the same altitude, making appearance-only classification
unreliable~\cite{benchmarking2024}. Third, the computational budget
is severely constrained: a model that achieves excellent accuracy on a
high-end GPU may be unusable on an edge device operating at 5--10
frames per second. Fourth, the problem extends beyond daylight hours
--- fires ignited at night or smouldering into the night are
effectively invisible to standard RGB models trained exclusively on
daytime data, and nighttime datasets remain
scarce~\cite{bouguettaya2022}. Finally, annotated datasets of
early-stage forest fire footage from fixed cameras in diverse geographic
regions remain limited, making generalisation to new deployment sites
such as South Asian forest terrain difficult~\cite{datasets20years}.

% -------------------------------------------------------
\section{Literature Review}
\label{sec:litreview}

This review organises the surveyed literature around the core challenges
facing the field rather than cataloguing papers individually.
Six challenges are identified: detection architecture design,
early-stage smoke and false-alarm suppression, cloud and smoke
confusion, nighttime and low-light detection, edge deployment and
model efficiency, and dataset availability and standardisation.
The emerging role of vision-language models is discussed last as a
cross-cutting theme.

% -------------------------------------------------------
\subsection{Detection Architectures: From Sensor Fusion to Deep CNNs}

Early work on automated fire perception relied on multi-sensor fusion
rather than data-driven learning. Martinez-de Dios et al.~\cite{martinez2008}
presented a system that combined visual and infrared (IR) cameras with
GPS telemetry to produce a real-time 3D model of a fire front, measuring
flame height, inclination angle, and base width simultaneously.
Fuzzy-wavelet filters were applied to suppress noise in the IR signal
while preserving the progressive spread of the fire front. Validated on
field experiments conducted in Gestosa, Portugal, between 1998 and 2005,
this system demonstrated that complementary information from visual and
infrared sensors substantially improves robustness over single-modality
approaches. However, the method is heavily hand-engineered, requiring
manual calibration and explicit feature design, and does not generalise
without significant re-engineering to new sites.

A follow-on line of work retained the IR modality but replaced
hand-crafted features with algorithmic segmentation. The approach
in~\cite{uav_ir2017} processes infrared UAV imagery by applying
histogram-based segmentation to extract hot candidate regions, then
uses optical flow to compute motion vectors for each candidate. Fire
regions are distinguished from fire analogues through morphological
operations and a blob counter. This pipeline achieves real-time fire
tracking within individual IR frames but requires a calibrated IR
sensor and is sensitive to abrupt camera movement.

The field shifted decisively toward deep learning following the success
of convolutional neural networks. The comprehensive review by
Bouguettaya et al.~\cite{bouguettaya2022} surveyed approximately 40
papers published between 2013 and 2021 and documented the transition
from hand-crafted RGB colour thresholds to end-to-end CNN architectures.
The dominant paradigm moved to one-stage detectors: multiple versions
of YOLO, SSD, and two-stage Faster R-CNN, as well as segmentation
networks such as U-Net and temporal models including LSTM and deep
belief networks. Subsequent work refined these architectures for the
fire domain. Kim and Muminov~\cite{kim2023} enhanced YOLOv7 with CBAM
attention, SPPF+, BiFPN, and decoupled heads, achieving AP50 of 86.4\%
on UAV smoke images. The evaluation in~\cite{yolov8_2024} further
reveals that YOLOv8 achieves recall of 0.9 for fire but only 0.7 for
smoke, quantifying the difficulty of early smoke detection
relative to visible flame.

% -------------------------------------------------------
\subsection{Early-Stage Smoke: Small Targets and False Alarms}

Detecting fire at the earliest possible stage means detecting smoke,
often before a visible flame exists. The Wildfire Smoke Detection
System in~\cite{wildfire2023} introduces the SKLFS-WildFire dataset,
which specifically captures fires in their initial stage with over
299,827 training images and proposes Cross Contrast Patch Embedding
(CCPE) to improve low-level smoke feature extraction in transformer
architectures. The authors demonstrate that standard transformers are
insufficient for this task because early-stage smoke texture and
transparency are encoded in low-level feature maps that transformer
patch embedding tends to discard. The Separable Negative Sampling
Mechanism (SNSM) additionally reduces false positives from background
clutter.

A continual learning variant of YOLOv3~\cite{fireecology2023}
incorporates a Feature Squeeze Block under a learning-without-forgetting
framework, achieving approximately 98\% detection precision while
retaining knowledge of previously seen deployment sites. This technique
is particularly relevant when a system must be updated with new
site-specific data without full re-training, as will be necessary when
in the case of new WWF camera sites or deployment in a new location. The comprehensive review
in~\cite{comprehensive2025} synthesises findings from 80 studies and
consistently identifies high false alarm rates as the primary unsolved
challenge across the literature.

% -------------------------------------------------------
\subsection{Cloud and Smoke Confusion in Fixed-Camera Deployments}

A particularly severe failure mode in fixed-camera setups is the
visual confusion between early-stage smoke and environmental phenomena
such as clouds, morning mist, steam from water bodies, and haze.
Unlike UAV systems that often look downward toward a ground reference,
PTZ cameras mounted at elevated forest positions look across the horizon,
placing clouds and distant smoke in the same spatial region of the frame
at the same apparent altitude. This makes appearance-only classification
unreliable at the most critical early-detection window --- dawn and dusk,
when mist and low cloud are most prevalent.

The benchmarking study in~\cite{benchmarking2024} explicitly addresses
this problem by constructing a multi-scene benchmark with hard negative
annotations for clouds, sunset conditions, morning mist, and steam.
The study demonstrates that models trained without these explicit hard
negatives produce false alarm rates 30--40\% higher than models trained
with them, confirming that the standard practice of simply omitting
confounders from training data is insufficient. The analysis
in~\cite{uncertainty2025} further shows that compact models
(YOLOv5n, YOLOv8n) are particularly susceptible to this failure mode
due to their limited representational capacity.

The cloud/smoke confusion problem operates at three distinct levels,
each requiring a different mitigation strategy. At the colour and
texture level, both clouds and smoke are grey-white with soft undefined
edges; explicit hard-negative training is the primary defence.
At the motion level, clouds drift uniformly across the frame while
smoke rises from a fixed ground point source and disperses irregularly;
temporal models that analyse consecutive frames can exploit this
motion difference~\cite{figlib2022}. At the contextual level,
smoke originates from a single location on the ground and its
trajectory follows wind direction, whereas clouds maintain consistent
altitude over longer time windows; scene-level context modelling
addresses this layer but remains largely unexplored. Critically,
this contextual understanding is precisely where vision-language
models such as CLIP offer a concrete advantage: by encoding descriptions
such as ``smoke rising from burning trees'' versus ``clouds drifting
over a mountain,'' CLIP provides a semantic disambiguation layer
that pure visual detectors fundamentally lack~\cite{clip_llava2026}.

% -------------------------------------------------------
\subsection{Nighttime and Low-Light Fire Detection}

Nighttime fire detection is one of the most underrepresented problems
in the field. The majority of published datasets and models assume
daylight conditions; Bouguettaya et al.~\cite{bouguettaya2022} note
this as an explicit gap, and the 20-year dataset review
in~\cite{datasets20years} confirms that night-time and low-light
annotated footage constitutes a small minority of all available data.
This is a practical problem: fires that ignite or smoulder at night
are invisible to RGB models trained exclusively on daytime imagery,
and the absence of colour cues means that the orange/red flame
signatures that dominate appearance-based detectors disappear entirely.

FLAME~2 is currently the most relevant publicly available dataset for
this sub-problem, providing seven paired RGB and infrared (IR) video
sequences captured during wildfire events. At night, IR sensors
detect heat signatures rather than visible light, making them
substantially more reliable than RGB cameras for flame localisation.
However, early-stage \textit{smoke} at night does not yet produce
a detectable heat signature, meaning IR alone is insufficient for
the earliest warning window.

A practical approach for projects with limited access to nighttime
data is synthetic augmentation: applying heavy darkening, blue
moonlight tinting, Gaussian sensor noise, and fire-region-preserving
darkness filters to existing daytime footage to create simulated
nighttime training samples. When combined with CLAHE contrast
enhancement (which approximates the behaviour of IR cameras by
amplifying local contrast differences), this approach can
substantially increase nighttime robustness without requiring a
dedicated nighttime collection campaign~\cite{kd2025}. The key
training strategy is staged fine-tuning: a model first trained on
large daytime datasets, then fine-tuned on available nighttime data
(FLAME~2 IR sequences), and finally adapted on site-specific footage
with synthetic nighttime augmentation applied to daytime WWF frames.

% -------------------------------------------------------
\subsection{Edge Deployment and Lightweight Model Design}

A critical constraint for this project is that inference must run in
real time on an edge device co-located with the PTZ camera.
Jadon et al.~\cite{firenet2019} were among the first to directly
target this constraint by designing FireNet, a custom CNN built from
scratch to balance detection accuracy against model size for
Raspberry Pi deployment. More recent work by~\cite{kd2025} applies
knowledge distillation to transfer accuracy from a powerful teacher
model to a compact MobileViT-S student, achieving strong detection
performance while meeting the computational budget of UAV-mounted
edge processors with concrete FPS benchmarks on embedded hardware.

An important practical note for this project is that edge device
performance benchmarks reported in the literature were obtained on
physical Jetson Nano or Raspberry Pi hardware. When direct access to
such devices is unavailable during development, as in the case of this project, laptop-based
simulation under constrained conditions (CPU thread capping, limited
VRAM allocation) provides a directional estimate but must not be
presented as equivalent to real hardware performance. Any reported
inference speed should explicitly state whether it was measured on
target hardware or simulated.

% -------------------------------------------------------
\subsection{Datasets and Benchmark Standardisation}

The quality and availability of annotated datasets is a bottleneck
across all sub-problems. The most extensive audit is provided
by~\cite{datasets20years}, which reviews datasets spanning 20 years.
FIgLib, introduced by Dewangan et al.~\cite{figlib2022}, is
particularly relevant to this project as it was collected from
fixed-view cameras analogous to PTZ setups. The dataset contains
315 fire sequences from 101 cameras in Southern California with
24,800 high-resolution images. The SmokeyNet architecture exploits
temporal structure through spatiotemporal feature aggregation,
confirming that consecutive-frame analysis substantially outperforms
single-frame detection for early smoke.

The benchmarking study in~\cite{benchmarking2024} identifies
inconsistent annotation standards across public datasets as a primary
source of unreliable cross-paper comparisons. FASDD provides over
120,000 heterogeneous images and is particularly useful for
generalisation testing because of its geographic and environmental
diversity. A notable limitation across all public datasets is
geographic bias: most are collected in North America or Europe and
do not reflect the vegetation, terrain, and atmospheric conditions
of South Asian forest environments targeted by this project,
reinforcing the necessity of a custom WWF dataset.

% -------------------------------------------------------
\subsection{Vision-Language Models for Contextual Fire Detection}

The most recent direction in the literature is the integration of
large vision-language models (VLMs) to address both the generalisation
and the false-alarm problems simultaneously.
Challa et al.~\cite{clip_llava2026} apply CLIP and LLaVA in a
zero-shot learning framework for forest fire detection, establishing
that pre-trained VLMs can classify fire/smoke images without any
domain-specific fine-tuning by leveraging natural language descriptions.
While zero-shot performance falls 10--20\% below fine-tuned supervised
methods, the paper demonstrates that combining CLIP features with a
fine-tuned detection head outperforms both approaches individually.

The most practically useful role of CLIP in this project's pipeline
is as a contextual disambiguation layer for low-confidence detections.
When a compact YOLO model produces a detection with confidence in the
0.3--0.6 range --- the ambiguous zone where cloud/smoke confusion is
most likely --- the detection crop can be passed to CLIP with targeted
text prompts: ``forest smoke rising from burning trees,''
``cloud drifting over mountain,'' ``morning mist in a valley.''
CLIP's vote, which encodes semantic and contextual understanding
beyond pixel-level appearance, then adjusts the final detection
confidence. This is especially powerful at dawn and dusk, and for
nighttime glow disambiguation (``fire glowing in dark forest'' versus
``vehicle headlight through trees''), without requiring nighttime
training data for CLIP itself.

\textbf{Open Problems (flagged by the field):}
\begin{itemize}[leftmargin=*, itemsep=2pt]
    \item False alarm rates from visual confounders (clouds, mist,
          sunset) remain the dominant unsolved challenge in deployed
          fixed-camera systems; no published method has achieved
          reliable performance under these conditions on held-out
          deployment sites~\cite{benchmarking2024, comprehensive2025}.
    \item Cloud and smoke confusion in horizon-facing fixed cameras
          is a distinct and largely unaddressed sub-problem; most
          datasets do not include sufficient hard negative examples
          of these confounders~\cite{benchmarking2024}.
    \item Nighttime and low-light fire detection lacks both adequate
          datasets and validated architectures; the field has
          identified it as a gap but no published solution has been
          demonstrated on fixed-camera forest
          footage~\cite{bouguettaya2022, datasets20years}.
    \item Domain shift from public datasets to new geographic
          deployments (e.g.\ South Asian forest terrain) is not
          solved by pre-processing alone and requires localised
          fine-tuning on site-specific footage, which is rarely
          studied~\cite{figlib2022}.
    \item Standardised evaluation benchmarks combining multiple
          datasets, consistent annotation formats, and realistic
          positive-to-negative ratios do not yet
          exist~\cite{datasets20years}.
    \item Efficient integration of temporal information from video
          on edge hardware has not been demonstrated at acceptable
          latency~\cite{figlib2022}.
\end{itemize}

% -------------------------------------------------------
\section{Research Gap and Proposed Direction}
\label{sec:gap}

The literature review exposes three compounding gaps. First, on the
accuracy side, high-performance models such as enhanced
YOLOv7~\cite{kim2023} and transformer-based
detectors~\cite{wildfire2023} require GPU-class hardware and have
not been validated on fixed PTZ camera footage in the geographic and
environmental conditions of this project. Second, on the deployment
side, lightweight models~\cite{firenet2019, kd2025} sacrifice
accuracy and are particularly unreliable for the harder sub-tasks
of early-stage smoke and visually ambiguous cloud/smoke
discrimination. Third, on the temporal and lighting side, nighttime
and low-light detection remains an open sub-problem with no published
solution validated on fixed-camera forest footage, and the interaction
between domain shift, staged fine-tuning, and synthetic nighttime
augmentation has not been studied in a single unified pipeline.

The WWF PTZ camera footage itself constitutes an additional gap:
no public dataset covers this geographic region, camera hardware,
or forest type. FIgLib is the closest analogue but is limited to
Southern California. Dataset construction and multi-condition
annotation (daytime, dawn/dusk, nighttime) from the raw PTZ footage
is therefore a necessary contribution in its own right.

Prior work achieves strong detection on curated benchmarks but
struggles with early-stage smoke, cloud/smoke confusion, and
nighttime conditions on edge hardware in real forest deployments;
we therefore propose to investigate a lightweight fire and smoke
detection model augmented with CLIP-based contextual features,
trained with staged fine-tuning and synthetic nighttime augmentation,
and evaluated on a custom annotated dataset from WWF-deployed PTZ
cameras under both daytime and low-light conditions.

\paragraph{Research Question (formal):}
Can a lightweight deep learning model, augmented with CLIP-based
contextual features and trained via staged fine-tuning with synthetic
nighttime augmentation, achieve real-time accurate early-stage fire
and smoke detection from fixed PTZ camera footage simulated on edge hardware, while outperforming standard compact-model
baselines on a custom WWF forest fire dataset?

% -------------------------------------------------------
\section{Benchmark, Evaluation, and Data Landscape}
\label{sec:benchmarks}

\paragraph{Datasets.}
Table~\ref{tab:datasets} summarises the key publicly available
datasets identified in the survey. FIgLib~\cite{figlib2022} and
SKLFS-WildFire~\cite{wildfire2023} are the most directly relevant
as primary training supplements. D-Fire serves as the standard
compact-model comparison benchmark~\cite{uncertainty2025}. FLAME~2
is the primary dataset for nighttime/IR training and FASDD provides
the largest heterogeneous collection for generalisation testing.

\begin{table}[h]
\caption{Key publicly available fire and smoke datasets.}
\label{tab:datasets}
\vskip 0.1in
\begin{center}
\begin{small}
\begin{tabular}{p{1.3cm} p{1.1cm} p{1.2cm} p{2.2cm}}
\toprule
\textbf{Dataset} & \textbf{Size} & \textbf{Modality}
& \textbf{Notes} \\
\midrule
D-Fire    & 21K+ img  & RGB       & Mixed outdoor; compact-model benchmark \\
FIgLib    & 24.8K img & RGB       & Fixed-view cameras; closest to PTZ setup \\
SKLFS     & 350K+ img & RGB       & Real early-stage wildfire footage \\
FLAME~2   & 7 videos  & RGB + IR  & Paired visual-IR; primary nighttime/low-light dataset \\
FASDD     & 120K+ img & RGB       & Heterogeneous; geographic diversity; good for generalisation \\
VisiFire  & 40 clips  & RGB       & Forest and urban scenes \\
FIRESENSE & 49 videos & RGB       & Heritage and urban environments \\
\bottomrule
\end{tabular}
\end{small}
\end{center}
\end{table}

\paragraph{Evaluation Protocol.}
The field uses mAP at IoU 0.5 (mAP\textsubscript{50}) as the primary
detection metric, with Precision, Recall, and F1-score alongside it.
For this project, Frames Per Second (FPS) on the target edge device
is an equally critical metric. False Positive Rate (FPR) under
confounder conditions (clouds, sunset, mist) must be reported
separately, as overall mAP masks the false-alarm problem that is the
dominant failure mode in deployment~\cite{benchmarking2024}.
Additionally, performance under simulated nighttime conditions
should be reported as a distinct evaluation partition.

\paragraph{Compute Context.}
Published methods primarily train on NVIDIA V100 or A100 GPUs.
Inference benchmarks in the edge deployment literature are conducted
on NVIDIA Jetson Nano (472 GFLOPS) and Raspberry Pi~4
($\approx$13.5 GFLOPS)~\cite{kd2025, firenet2019}. Quantised and
distilled models typically achieve 10--30~FPS on Jetson Nano;
unoptimised YOLOv8 variants achieve approximately 3--8~FPS on the
same hardware.

% -------------------------------------------------------
\section*{Experimental Design Considerations}
\label{sec:expdesign}

\begin{itemize}[leftmargin=*, itemsep=2pt]
    \item \textbf{Baselines.} Any proposed method must outperform at
          minimum YOLOv8n~\cite{uncertainty2025} and
          FireNet~\cite{firenet2019} on both mAP\textsubscript{50}
          and FPS on target edge hardware. A CLIP zero-shot
          baseline~\cite{clip_llava2026} must also be evaluated to
          quantify the benefit of fine-tuning over zero-shot
          inference. A daytime-only model trained without nighttime
          augmentation should be evaluated as a separate ablation
          point to isolate the contribution of nighttime training.

    \item \textbf{Ablation.} Key variables to isolate: (1) backbone
          size vs.\ accuracy trade-off; (2) marginal contribution of
          CLIP features over a purely visual backbone; (3) effect of
          including cloud/mist/sunset hard negatives on false alarm
          rate; (4) contribution of synthetic nighttime augmentation
          on low-light detection performance.

    \item \textbf{Evaluation pitfalls.}
          \textit{Geographic leakage:} training and testing on the
          same camera site inflates generalisation estimates; at
          least one WWF camera site must be held out
          entirely~\cite{figlib2022}.
          \textit{Positive-negative imbalance:} the 50/50 split in
          FIgLib does not reflect real deployment ratios; precision
          at low false alarm rates is more informative than mAP
          alone.
          \textit{Confounder blindness:} clouds, sunset, and
          morning mist must be tested as an explicit evaluation
          partition; models routinely fail on these even when
          overall mAP is acceptable~\cite{benchmarking2024}.
          \textit{Simulation vs.\ real hardware:} if Jetson or
          Raspberry Pi devices are unavailable, inference speed must
          be measured under CPU-constrained conditions and explicitly
          reported as simulated; GPU benchmark numbers must not be
          presented as edge-device performance.
\end{itemize}

% -------------------------------------------------------
\section{Annotated Reading List}
\label{sec:readings}

\begin{thebibliography}{99}

\bibitem{martinez2008}
J.~R.~Martinez-de Dios, B.~C.~Arrue, A.~Ollero, L.~Merino, and
F.~G\'{o}mez-Rodr\'{i}guez.
\newblock Computer vision techniques for forest fire perception.
\newblock \textit{Image and Vision Computing}, 26(5):550--562, 2008.
\newblock \textit{What it does: Fuses visual camera, IR camera, and GPS
to build a real-time 3D fire model using fuzzy-wavelet filters.}
\newblock \textit{What it contributes: Demonstrates IR + RGB sensor
fusion improves measurement accuracy over single-modality approaches.}
\newblock \textit{Why it is here: Foundational sensor-fusion baseline;
motivates multi-modal approaches and IR use for nighttime detection.}

\bibitem{bouguettaya2022}
A.~Bouguettaya, H.~Zarzour, A.~M.~Taberkit, and A.~Kechida.
\newblock A review on early wildfire detection from UAVs using deep
learning-based computer vision algorithms.
\newblock \textit{Signal Processing}, 190:108309, 2022.
\newblock \textit{What it does: Surveys $\sim$40 papers from 2013--2021
covering YOLO, R-CNN, SSD, U-Net, LSTM, and communication technologies.}
\newblock \textit{What it contributes: Most comprehensive landscape
overview; identifies nighttime detection and edge complexity as key gaps.}
\newblock \textit{Why it is here: Primary landscape paper covering all
architectures we compare against.}

\bibitem{uav_ir2017}
(Authors confirmed via IEEE Xplore arXiv no.~7991306.)
\newblock Fire detection using infrared images for UAV-based forest fire
surveillance.
\newblock \textit{IEEE Conference Proceedings}, 2017.
\newblock \textit{What it does: Detects fire in UAV IR imagery via
histogram segmentation and optical flow motion analysis.}
\newblock \textit{What it contributes: Shows brightness and motion in IR
provides reliable fire localisation without deep learning.}
\newblock \textit{Why it is here: Classical IR baseline; motivates IR
as a complement to RGB for nighttime fire localisation.}

\bibitem{firenet2019}
A.~Jadon, M.~Omama, A.~Varshney, M.~S.~Ansari, and R.~Sharma.
\newblock {FireNet}: A specialized lightweight fire \& smoke detection
model for real-time {IoT} applications.
\newblock \textit{arXiv:1905.11922}, 2019.
\newblock \textit{What it does: Designs a shallow CNN (FireNet) from
scratch for Raspberry Pi real-time fire/smoke detection.}
\newblock \textit{What it contributes: Establishes the lightweight CNN
edge baseline; open-sources model and dataset.}
\newblock \textit{Why it is here: Primary compact-model baseline for
Raspberry Pi / Jetson deployment comparison.}

\bibitem{kim2023}
S.-Y.~Kim and A.~Muminov.
\newblock Forest fire smoke detection based on deep learning approaches
and unmanned aerial vehicle images.
\newblock \textit{Sensors}, 23(12), 2023.
\newblock \textit{What it does: Enhances YOLOv7 with CBAM, SPPF+,
BiFPN, and decoupled heads on 6,500 UAV smoke images.}
\newblock \textit{What it contributes: AP50 of 86.4\%; best published
UAV smoke detection result in this survey.}
\newblock \textit{Why it is here: Upper-bound accuracy target before
edge optimisation.}

\bibitem{yolov8_2024}
(Authors from IJACSA Vol.~15 No.~1.)
\newblock A {YOLO}-based approach for fire and smoke detection.
\newblock \textit{IJACSA}, 15(1), 2024.
\newblock \textit{What it does: Applies YOLOv8 and separately analyses
precision-recall curves for fire vs.\ smoke classes.}
\newblock \textit{What it contributes: Quantifies the fire/smoke
performance gap (recall 0.9 vs.\ 0.7 at same threshold).}
\newblock \textit{Why it is here: Provides YOLOv8 baseline and
motivates special treatment for the smoke sub-task.}

\bibitem{wildfire2023}
(Authors from arXiv:2311.10116.)
\newblock Wildfire smoke detection system: Model architecture, training
mechanism, and dataset.
\newblock \textit{arXiv:2311.10116}, 2023.
\newblock \textit{What it does: Introduces SKLFS-WildFire (350K+ early
wildfire images) and CCPE + SNSM modules on YOLOX.}
\newblock \textit{What it contributes: Largest real early-stage wildfire
dataset; shows transformers alone are insufficient for early smoke.}
\newblock \textit{Why it is here: Hardest sub-problem in our project;
SKLFS supplements WWF data during training.}

\bibitem{fireecology2023}
(Authors from Fire Ecology, Springer Open.)
\newblock Forest fire and smoke detection using deep learning-based
learning without forgetting.
\newblock \textit{Fire Ecology}, 2022/2023.
\newblock \textit{What it does: YOLOv3 + Feature Squeeze Block under
continual learning; $\sim$98\% precision retaining prior site knowledge.}
\newblock \textit{What it contributes: Demonstrates LwF prevents
catastrophic forgetting when fine-tuning on new site data.}
\newblock \textit{Why it is here: Relevant when the model must be
updated as new WWF camera sites come online.}

\bibitem{comprehensive2025}
(Authors from MDPI Applied Sciences, 2025.)
\newblock Early fire and smoke detection using deep learning: A
comprehensive review.
\newblock \textit{Applied Sciences}, 15(18):10255, 2025.
\newblock \textit{What it does: Reviews 80 studies across CNNs, YOLO,
Faster R-CNN, and spatiotemporal models.}
\newblock \textit{What it contributes: Identifies high false alarm rates
as the dominant unsolved problem across the field.}
\newblock \textit{Why it is here: Cross-validates open problems identified
in this survey; ensures complete coverage.}

\bibitem{benchmarking2024}
(Authors from arXiv:2410.16631.)
\newblock Benchmarking multi-scene fire and smoke detection.
\newblock \textit{arXiv:2410.16631}, 2024.
\newblock \textit{What it does: Unifies FIRESENSE, BoWFire, D-Fire with
consistent labels and explicit confounder annotations.}
\newblock \textit{What it contributes: Shows confounder training reduces
false alarms 30--40\%; exposes annotation inconsistency across datasets.}
\newblock \textit{Why it is here: Directly informs our cloud/smoke hard
negative strategy and evaluation protocol.}

\bibitem{figlib2022}
A.~Dewangan, Y.~Pande, H.-W.~Braun, F.~Vernon, I.~Perez, I.~Altintas,
G.~W.~Cottrell, and M.~H.~Nguyen.
\newblock {FIgLib} \& {SmokeyNet}: Dataset and deep learning model for
real-time wildland fire smoke detection.
\newblock \textit{Remote Sensing}, 14(4):1007, 2022.
\newblock \textit{What it does: Introduces FIgLib (24,800 images from
101 fixed-view cameras) and SmokeyNet spatiotemporal architecture.}
\newblock \textit{What it contributes: Closest PTZ-camera analogue;
temporal model rivals human performance in smoke detection.}
\newblock \textit{Why it is here: Primary pre-training dataset;
SmokeyNet is key architectural reference for temporal detection.}

\bibitem{datasets20years}
(Authors from arXiv:2503.14552.)
\newblock Fire and smoke datasets in 20 years: An in-depth review.
\newblock \textit{arXiv:2503.14552}, 2025.
\newblock \textit{What it does: Reviews every major fire/smoke dataset
from 20 years across RGB, thermal, IR; benchmarks on ResNet-50,
DeepLabV3, YOLOv8.}
\newblock \textit{What it contributes: Exposes cross-dataset performance
variation; most comprehensive dataset audit available.}
\newblock \textit{Why it is here: Essential for dataset selection and
understanding generalisation limits to new geographies.}

\bibitem{kd2025}
(Authors from arXiv:2502.20979.)
\newblock Real-time aerial fire detection on resource-constrained devices
using knowledge distillation.
\newblock \textit{arXiv:2502.20979}, 2025.
\newblock \textit{What it does: Compresses MobileViT-S via knowledge
distillation for UAV edge processors with hardware FPS benchmarks.}
\newblock \textit{What it contributes: State-of-the-art edge compression;
concrete latency numbers on embedded hardware.}
\newblock \textit{Why it is here: Edge compression technique to apply
or compare against for Jetson / Raspberry Pi deployment.}

\bibitem{uncertainty2025}
A.~S.~Joshi et al.
\newblock Uncertainty-aware post-detection framework for enhanced fire
and smoke detection in compact deep learning models.
\newblock \textit{arXiv:2510.10108}, 2025.
\newblock \textit{What it does: Evaluates YOLOv5n/YOLOv8n on D-Fire;
proposes Confidence Refinement Network using uncertainty + colour/edge.}
\newblock \textit{What it contributes: Lightweight post-detection module
improving compact-model performance without modifying the base model.}
\newblock \textit{Why it is here: Actionable technique for improving
compact baseline; benchmarks on D-Fire our standard eval dataset.}

\bibitem{clip_llava2026}
N.~R.~Challa, B.~C.~Mohan, and M.~L.~Babu.
\newblock Enhancing forest fire detection using zero-shot learning:
A {LLaVA} and {CLIP} vision-language model approach.
\newblock \textit{Multimedia Tools and Applications}, Springer, 2026.
\newblock \textit{What it does: Applies CLIP and LLaVA zero-shot for
forest fire detection using natural language prompts.}
\newblock \textit{What it contributes: First zero-shot VLM baseline for
forest fire; combining CLIP with fine-tuned head outperforms both alone.}
\newblock \textit{Why it is here: Motivates and benchmarks CLIP in our
pipeline; defines zero-shot lower bound and cloud/smoke disambiguation
use case.}

\end{thebibliography}

% -------------------------------------------------------
\section{Preliminary Implementation Resources}
\label{sec:resources}

\begin{itemize}[leftmargin=*, itemsep=2pt]
    \item \url{https://github.com/ultralytics/ultralytics}
          YOLOv8/v11 repo. Primary backbone candidate with TensorRT
          export for Jetson deployment.

    \item \url{https://github.com/ultralytics/yolov5}
          YOLOv5 nano/small variants. Training infrastructure for
          compact-model baselines.

    \item \url{http://hpwren.ucsd.edu/HPWREN-FIgLib/}
          FIgLib dataset download. Primary fixed-camera pre-training
          dataset; closest public analogue to PTZ setup.

    \item \url{https://github.com/openai/CLIP}
          OpenAI CLIP (ViT-B/32). Visual-language embeddings for
          contextual cloud/smoke disambiguation and nighttime glow
          classification.

    \item \url{https://github.com/facebookresearch/detectron2}
          Facebook Detectron2. Faster R-CNN / Mask R-CNN baseline
          framework; referenced in the project description.

    \item \url{https://universe.roboflow.com}
          Roboflow dataset hub. YOLO-format fire/smoke datasets for
          augmenting WWF annotations; also provides nighttime fire
          image collections.

    \item \url{https://github.com/AnshP99/firenet}
          FireNet open-source implementation. Lightweight edge
          baseline to reproduce and benchmark on target hardware.

    \item \url{https://arxiv.org/abs/2112.08598}
          FIgLib \& SmokeyNet arXiv preprint. Full paper, dataset
          documentation, and SmokeyNet architecture details.
\end{itemize}

% -------------------------------------------------------
% PAPER LINKS AND GITHUB RESOURCES
% -------------------------------------------------------
\section*{Full Paper Links and Resources}
\label{sec:links}

\subsection*{Research Paper Links}

\begin{enumerate}[leftmargin=*, itemsep=4pt]

    \item \textbf{[P1] Computer Vision Techniques for Forest Fire
          Perception} --- Martinez-de Dios et al., 2008\\
          \url{https://www.sciencedirect.com/science/article/pii/S0262885607001096}

    \item \textbf{[P2] A Review on Early Wildfire Detection from UAVs
          Using Deep Learning} --- Bouguettaya et al., 2022\\
          \url{https://www.sciencedirect.com/science/article/pii/S0165168421003467}

    \item \textbf{[P3] Fire Detection Using Infrared Images for
          UAV-Based Forest Fire Surveillance} --- IEEE, 2017\\
          \url{https://ieeexplore.ieee.org/document/7991306}

    \item \textbf{[P4] FireNet: A Specialized Lightweight Fire \&
          Smoke Detection Model} --- Jadon et al., 2019\\
          \url{https://arxiv.org/abs/1905.11922}

    \item \textbf{[P5] Forest Fire Smoke Detection Based on Deep
          Learning and UAV Images} --- Kim \& Muminov, 2023\\
          \url{https://pmc.ncbi.nlm.nih.gov/articles/PMC10304711/}

    \item \textbf{[P6] A YOLO-Based Approach for Fire and Smoke
          Detection} --- IJACSA, 2024\\
          \url{https://thesai.org/Downloads/Volume15No1/Paper_9-A_Yolo_based_Approach_for_Fire_and_Smoke_Detection.pdf}

    \item \textbf{[P7] Wildfire Smoke Detection System: Model
          Architecture, Training, and Dataset} --- arXiv, 2023\\
          \url{https://arxiv.org/abs/2311.10116}

    \item \textbf{[P8] Forest Fire and Smoke Detection Using Learning
          Without Forgetting} --- Fire Ecology, 2023\\
          \url{https://link.springer.com/article/10.1186/s42408-022-00165-0}

    \item \textbf{[P9] Early Fire and Smoke Detection: A Comprehensive
          Review} --- MDPI Applied Sciences, 2025\\
          \url{https://www.mdpi.com/2076-3417/15/18/10255}

    \item \textbf{[P10] Benchmarking Multi-Scene Fire and Smoke
          Detection} --- arXiv, 2024\\
          \url{https://arxiv.org/abs/2410.16631}

    \item \textbf{[P11] FIgLib \& SmokeyNet: Dataset and Model for
          Wildland Fire Smoke Detection} --- Remote Sensing, 2022\\
          \url{https://arxiv.org/abs/2112.08598}\\
          Dataset: \url{http://hpwren.ucsd.edu/HPWREN-FIgLib/}

    \item \textbf{[P12] Fire and Smoke Datasets in 20 Years:
          An In-Depth Review} --- arXiv, 2025\\
          \url{https://arxiv.org/abs/2503.14552}

    \item \textbf{[P13] Real-Time Aerial Fire Detection Using
          Knowledge Distillation} --- arXiv, 2025\\
          \url{https://arxiv.org/abs/2502.20979}

    \item \textbf{[P14] Uncertainty-Aware Post-Detection Framework
          for Fire and Smoke Detection} --- arXiv, 2025\\
          \url{https://arxiv.org/abs/2510.10108}

    \item \textbf{[P15] Enhancing Forest Fire Detection Using
          Zero-Shot Learning: CLIP and LLaVA} --- Springer, 2026\\
          \url{https://link.springer.com/article/10.1007/s11042-026-21377-5}

\end{enumerate}

\subsection*{GitHub Repositories}

\begin{enumerate}[leftmargin=*, itemsep=4pt]

    \item \textbf{Ultralytics YOLOv8/v11} --- Primary detection backbone
          with TensorRT / ONNX export for Jetson Nano.\\
          \url{https://github.com/ultralytics/ultralytics}

    \item \textbf{Ultralytics YOLOv5} --- Compact nano/small variants
          for edge-baseline experiments.\\
          \url{https://github.com/ultralytics/yolov5}

    \item \textbf{OpenAI CLIP} --- ViT-B/32 visual-language embeddings
          for cloud/smoke disambiguation and nighttime glow
          classification.\\
          \url{https://github.com/openai/CLIP}

    \item \textbf{Facebook Detectron2} --- Faster R-CNN / Mask R-CNN
          framework; referenced baseline in the project description.\\
          \url{https://github.com/facebookresearch/detectron2}

    \item \textbf{FireNet (AnshP99)} --- Open-source FireNet with
          training dataset; lightweight edge baseline to reproduce.\\
          \url{https://github.com/AnshP99/firenet}

    \item \textbf{Roboflow Universe} --- YOLO-format fire/smoke datasets
          including nighttime collections for supplementing WWF data.\\
          \url{https://universe.roboflow.com/}

\end{enumerate}

\end{document}
